
\documentclass{article}
\usepackage{colortbl}
\usepackage{makecell}
\usepackage{multirow}
\usepackage{supertabular}

\begin{document}

\newcounter{utterance}

\centering \large Interaction Transcript for game `eqbench', experiment `eq\_bench', episode 79 with reflexion{-}llama{-}3.1{-}8b{-}instruct.
\vspace{24pt}

{ \footnotesize  \setcounter{utterance}{1}
\setlength{\tabcolsep}{0pt}
\begin{supertabular}{c@{$\;$}|p{.15\linewidth}@{}p{.15\linewidth}p{.15\linewidth}p{.15\linewidth}p{.15\linewidth}p{.15\linewidth}}
    \# & \multicolumn{2}{c}{Player} && \multicolumn{2}{c}{Game Master} \\
    \hline

    \theutterance \stepcounter{utterance}  
    & & & \multicolumn{4}{p{0.6\linewidth}}{
        \cellcolor[rgb]{0.9,0.9,0.9}{
            \makecell[{{p{\linewidth}}}]{
                \texttt{\tiny{[P1$\langle$GM]}}
                \texttt{You have attempted this same task before and did not succeed. These are the plans you wrote after those attempts.} \\
\texttt{Trial \#1: New plan: I got stuck in a loop of assigning high scores to emotions that are not likely to be felt by Bennett in this situation. Upon re{-}reading the dialogue, I realize that Bennett is trying to reason with Julian and provide evidence to counter his conspiracy theory. Bennett's tone is more likely to be defensive and irritated, as he is being challenged by Julian's claims. However, Bennett is not feeling amusement, as the conversation is serious and not humorous. Bennett is also not feeling relief, as the conversation is not resolving or ending in a positive way. Next time, I will assign scores based on the context and tone of the dialogue, taking into account the character's actions and words. I will also consider the possibility that Bennett may be feeling a mix of emotions, such as a combination of irritation and defensiveness.} \\
\\ 
\\ 
\\ 
\texttt{Your task is to predict the likely emotional responses of a character in this dialogue:} \\
\\ 
\texttt{Bennett: Julian, I've been in space. I've seen the Earth from a distance. It's round. } \\
\texttt{Julian: That's what they want you to believe, Bennett.} \\
\texttt{Bennett: And who exactly are 'they'?} \\
\texttt{Julian: The ones pulling the strings, the ones in power. } \\
\texttt{Bennett: You mean scientists? People dedicated to exploring and understanding our universe? } \\
\texttt{Julian: No, the puppeteers. The ones who manipulate the scientists. } \\
\texttt{Bennett: You're suggesting a global conspiracy, Julian. That's a serious claim. } \\
\texttt{Julian: And I stand by it. } \\
\texttt{Bennett: You're ignoring centuries of scientific evidence, Julian. } \\
\texttt{Julian: I'm not ignoring it, Bennett. I'm questioning it. } \\
\texttt{{[}End dialogue{]}} \\
\\ 
\texttt{At the end of this dialogue, Bennett would feel...} \\
\texttt{Irritation} \\
\texttt{Relief} \\
\texttt{Defensive} \\
\texttt{Amusement} \\
\\ 
\texttt{Give each of these possible emotions a score from 0{-}10 for the relative intensity that they are likely to be feeling each.} \\
\\ 
\texttt{You must output in the following format, including headings (of course, you should give your own scores), with no additional commentary:} \\
\\ 
\texttt{Irritation: <score>} \\
\texttt{Relief: <score>} \\
\texttt{Defensive: <score>} \\
\texttt{Amusement: <score>} \\
\\ 
\\ 
\texttt{{[}End of answer{]}} \\
\\ 
\texttt{Remember: zero is a valid score, meaning they are likely not feeling that emotion. You must score at least one emotion > 0.} \\
\\ 
\texttt{Your answer:} \\
            }
        }
    }
     \\ \\

    \theutterance \stepcounter{utterance}  
    & \multicolumn{4}{p{0.6\linewidth}}{
        \cellcolor[rgb]{0.561,0.737,0.561}{
            \makecell[{{p{\linewidth}}}]{
                \texttt{\tiny{[P1$\rangle$GM]}}
                \texttt{Irritation: 8} \\
\texttt{Relief: 0} \\
\texttt{Defensive: 9} \\
\texttt{Amusement: 0} \\
            }
        }
    }
    & & \\ \\

    \theutterance \stepcounter{utterance}  
    & & & \multicolumn{2}{p{0.3\linewidth}}{
        \cellcolor[rgb]{0.9,0.9,0.9}{
            \makecell[{{p{\linewidth}}}]{
                \texttt{\tiny{[GM$|$GM]}}
                \texttt{Irritation: 8} \\
\texttt{Relief: 0} \\
\texttt{Defensive: 9} \\
\texttt{Amusement: 0} \\
            }
        }
    }
    & & \\ \\

    \theutterance \stepcounter{utterance}  
    & & & \multicolumn{2}{p{0.3\linewidth}}{
        \cellcolor[rgb]{0.9,0.9,0.9}{
            \makecell[{{p{\linewidth}}}]{
                \texttt{\tiny{[GM$|$GM]}}
                \texttt{Irritation: 8} \\
\texttt{Relief: 0} \\
\texttt{Defensive: 3} \\
\texttt{Amusement: 2} \\
            }
        }
    }
    & & \\ \\

    \theutterance \stepcounter{utterance}  
    & & & \multicolumn{2}{p{0.3\linewidth}}{
        \cellcolor[rgb]{0.9,0.9,0.9}{
            \makecell[{{p{\linewidth}}}]{
                \texttt{\tiny{[GM$|$GM]}}
                \texttt{game\_result = LOSE} \\
            }
        }
    }
    & & \\ \\

    \theutterance \stepcounter{utterance}  
    & & \multicolumn{4}{p{0.6\linewidth}}{
        \cellcolor[rgb]{0.95,0.95,0.95}{
            \makecell[{{p{\linewidth}}}]{
                \texttt{\tiny{[P1$\leftrightarrow$P]}}
                \texttt{Post{-}episode Reflexion (not sent to GM):} \\
\texttt{New plan: Upon re{-}reading the dialogue, I realize that I overestimated the intensity of Bennett's emotions, particularly amusement and relief. I also failed to consider the possibility that Bennett may be feeling a mix of emotions. Next time, I will take a more nuanced approach by analyzing the dialogue and considering the context, tone, and character's actions and words. I will assign scores based on the likelihood of each emotion being felt by Bennett, taking into account the conversation's progression and the character's reactions. I will also consider the possibility of Bennett feeling a combination of emotions, such as irritation and defensiveness. To achieve this, I will:} \\
\\ 
\texttt{1. Identify the key points in the dialogue where Bennett's emotions are likely to be triggered.} \\
\texttt{2. Analyze Bennett's tone and language to determine the intensity of his emotions.} \\
\texttt{3. Consider the context of the conversation and how it may influence Bennett's emotional state.} \\
\texttt{4. Assign scores based on the likelihood of each emotion being felt by Bennett, taking into account the conversation's progression and the character's reactions.} \\
\\ 
\texttt{By taking a more nuanced approach, I aim to provide a more accurate assessment of Bennett's emotional state and assign scores that reflect the complexity of human emotions.} \\
            }
        }
    }
    & \\ \\

\end{supertabular}
}

\end{document}
